%============== Bảng màu VinhMath đồng bộ với nhận diện thương hiệu Website
% Load xcolor package trước khi định nghĩa màu
\usepackage[dvipsnames,table]{xcolor}

% --- MÀU SẮC THƯƠNG HIỆU VINHMATH (HỆ ĐƠN SẮC VÀNG) ---
\definecolor{vmAccent}{HTML}{854D0E}     % Vàng đất đậm (thay xanh Navy) - Điểm nhấn chính
\definecolor{vmAccentSoft}{HTML}{FEF3C7} % Nền vàng nhạt / vàng cát
\definecolor{vmOk}{HTML}{C27D00}         % Vàng trung tính
\definecolor{vmOkSoft}{HTML}{FFFDF0}     % Nền vàng nhạt
\definecolor{vmWarn}{HTML}{D97706}       % Vàng tươi sáng (thay màu Đỏ)
\definecolor{vmWarnSoft}{HTML}{FFFBEB}   % Nền vàng tươi nhạt
\definecolor{vmInk}{HTML}{1C1917}        % Mực Charcoal đen ấm (trên nền trắng)
\definecolor{vmBg}{HTML}{FFFFFF}         % Nền giấy trắng tinh khiết
\definecolor{vmLine}{HTML}{E7E5E4}       % Đường kẻ xám ấm
\definecolor{warmgray}{HTML}{78716C}     % Xám ấm chú thích

% --- Đè màu chuẩn trong TikZ/LaTeX để đồng bộ thương hiệu VinhMath ---
\colorlet{blue}{vmAccent}      % Nét vẽ chính -> Vàng đất đậm (phong cách Sepia)
\colorlet{green}{warmgray}     % Nét vẽ phụ/chú thích -> Xám ấm
\colorlet{magenta}{vmOk}       % Cung tròn góc thường -> Vàng trung tính
\colorlet{orange}{vmWarn}      % Góc cần tính/điểm nhấn -> Vàng tươi sáng
\colorlet{violet}{vmAccent}
\colorlet{cyan}{vmAccent}
\colorlet{teal}{vmAccent}

% --- Mapping tương thích ngược với phong cách cũ ---
\definecolor{acadbg}{HTML}{FFFFFF}       % Nền trắng tinh cho giấy in
\definecolor{acadfg}{HTML}{1F1D1A}       % Chữ đen Charcoal
\definecolor{acadgray}{HTML}{64748B}     % Xám chú thích
\definecolor{acadlightgray}{HTML}{F8FAFC}% Xám rất nhạt
\definecolor{acadrule}{HTML}{E2E8F0}     % Đường kẻ mảnh
\definecolor{acadnavy}{HTML}{0F2942}     % Accent Xanh dương đen chính
\definecolor{acadnavylight}{HTML}{C27D00}% Accent Vàng đậm phụ

% Mapping cho code cũ (Dracula & Minimal cũ)
\definecolor{dracbg}{HTML}{FFFFFF}
\definecolor{dracline}{HTML}{E2E8F0}
\definecolor{dracfg}{HTML}{1F1D1A}
\definecolor{draccomment}{HTML}{64748B}
\definecolor{draccyan}{HTML}{0F2942}
\definecolor{dracgreen}{HTML}{C27D00}
\definecolor{dracorange}{HTML}{B5483A}
\definecolor{dracpink}{HTML}{0F2942}
\definecolor{dracpurple}{HTML}{0F2942}
\definecolor{dracred}{HTML}{B5483A}
\definecolor{dracyellow}{HTML}{C27D00}

\definecolor{acateal}{HTML}{0F2942}
\definecolor{acadarkteal}{HTML}{0F2942}
\definecolor{acalightteal}{HTML}{EFF6FF}
\definecolor{acalightcyan}{HTML}{FFFDF0}
\definecolor{acagray}{HTML}{1F1D1A}
\definecolor{acadarkgray}{HTML}{0F2942}
\definecolor{acalightgray}{HTML}{FFFFFF}
\definecolor{acaborderblue}{HTML}{0F2942}
\definecolor{acadeepblue}{HTML}{0F2942}
\definecolor{acarust}{HTML}{64748B}

% Màu trang bìa (giữ tương thích và đổi sang tông VinhMath)
\definecolor{ivory}{HTML}{FFFFFF}
\definecolor{lapislazuli}{HTML}{0F2942}
\definecolor{terracotta}{HTML}{C27D00}
\definecolor{darkslate}{HTML}{1F1D1A}
\definecolor{oldgold}{HTML}{B5483A}
\definecolor{marblegray}{HTML}{E2E8F0}

% --- Cấu hình màu cho môi trường toán học ---
\def\mauVD{vmAccent}                      % Ví dụ dùng Đỏ son Hanko
\def\mauEX{vmOk}                          % Bài tập ví dụ dùng Xanh Matcha
\def\mauBT{vmWarn}                        % Bài tập tự luyện dùng Vàng hổ phách
\def\mauLG{warmgray}                      % Lời giải & đường chấm dùng Xám ấm
\def\mauDA{\mauEX}
\def\mauTrue{vmAccent}

%%%%%% Màu mục lục
\def\mauPART{vmAccent}
\def\mauCHAP{vmAccent}
\def\mauSEC{vmOk}
\def\mauSUBSEC{vmWarn}
\def\mauSUBSUBSEC{vmInk}

%%%%%% Màu các tiêu đề
\def\mauphan{vmAccent}
\def\mauchuphan{vmAccent}
\def\mauchuong{vmAccent}
\def\maunenmucluc{white}
\def\maumlmot{vmAccent}
\def\maumlhai{vmOk}
\def\maumlba{vmWarn}
\def\maubai{vmOk}
\def\mausubsec{vmWarn}
\def\mausubsecphu{vmAccent}
\def\mausubsubsec{vmInk}
\def\maupara{vmInk}

%%%%%% Footer/Header
\def\maufoot{warmgray}
\def\mauchufoot{warmgray}
\def\maufootmot{warmgray}
\def\maufoothai{warmgray}
\def\maufootba{warmgray}

%%%%%% Dạng toán
\def\maudang{vmAccent}
\def\maunendang{white}

%%%%%% Định nghĩa, định lý (dùng tcolorbox)
\def\maudl{vmInk}
\def\nendl{vmBg}
\def\maudn{vmAccent}
\def\nendn{vmAccentSoft}

\def\maucy{vmAccent}
\def\maupt{vmInk}

%%%%%% Dấu tích đầu câu
\def\mauitemCI{vmAccent}
\def\mauitemKN{vmOk}
\def\itemKN{\color{\mauitemKN}\faCheckSquareO}
\def\itemCI{\color{\mauitemCI}\faCheckCircleO}
